\documentclass[11pt,a4paper]{article}
\usepackage[utf8]{inputenc}
\usepackage[spanish]{babel}
\usepackage[margin=2.5cm]{geometry}
\usepackage{xcolor}
\usepackage{listings}
\usepackage{tcolorbox}
\usepackage{tikz}
\usepackage{hyperref}
\usepackage{enumitem}

% Configuración de colores
\definecolor{codegreen}{rgb}{0,0.6,0}
\definecolor{codegray}{rgb}{0.5,0.5,0.5}
\definecolor{codepurple}{rgb}{0.58,0,0.82}
\definecolor{backcolour}{rgb}{0.95,0.95,0.92}

% Configuración de listings
\lstdefinestyle{mystyle}{
    backgroundcolor=\color{backcolour},
    commentstyle=\color{codegreen},
    keywordstyle=\color{magenta},
    numberstyle=\tiny\color{codegray},
    stringstyle=\color{codepurple},
    basicstyle=\ttfamily\small,
    breakatwhitespace=false,
    breaklines=true,
    captionpos=b,
    keepspaces=true,
    numbers=left,
    numbersep=5pt,
    showspaces=false,
    showstringspaces=false,
    showtabs=false,
    tabsize=2,
    frame=single
}
\lstset{style=mystyle}

% Configuración de tcolorbox
\tcbuselibrary{skins,breakable}
\newtcolorbox{notebox}[1][]{
    colback=blue!5!white,
    colframe=blue!75!black,
    fonttitle=\bfseries,
    title=#1
}

\title{\textbf{Guía de Esteganografía de Doble Capa}\\
\large Uso de Steghide y OpenSSL para Comunicación Segura}
\author{}
\date{\today}

\begin{document}

\maketitle

\tableofcontents
\newpage

\section{Concepto General}

Este método implementa un esquema de \textbf{doble capa de seguridad} utilizando esteganografía y cifrado simétrico. La idea es ocultar información en dos imágenes separadas:

\begin{table}[h]
\centering
\begin{tabular}{|l|l|l|}
\hline
\textbf{Imagen} & \textbf{Contenido} & \textbf{Función} \\
\hline
\textbf{RedTeam.jpg} & Texto plano (la "llave") & Es la contraseña \\
\hline
\textbf{BlueTeam.jpg} & Archivo cifrado con esa llave & Es el cofre \\
\hline
\end{tabular}
\end{table}

El receptor \textbf{primero} extrae la llave de RedTeam, y \textbf{después} la usa para abrir el cofre de BlueTeam.

\section{Fase 1: Creación (Emisor)}

\subsection{Paso 1 --- Definir la llave y el mensaje secreto}

\begin{lstlisting}[language=bash, caption={Crear archivos iniciales}]
# La llave que ocultarás en RedTeam (puede ser cualquier texto)
echo "CaballoDeTroya2026!" > RedKey.txt

# El mensaje secreto real que quieres que lea el receptor
echo "El punto de encuentro es el muelle 7 a las 03:00 AM. Destruye este mensaje." > mensaje_secreto.txt
\end{lstlisting}

\subsection{Paso 2 --- Cifrar el mensaje secreto usando la llave}

Usamos \texttt{openssl} para cifrar \texttt{mensaje\_secreto.txt} con la contraseña que está dentro de \texttt{RedKey.txt}:

\begin{lstlisting}[language=bash, caption={Cifrar mensaje con OpenSSL}]
openssl enc -aes-256-cbc -salt -pbkdf2 \
  -in mensaje_secreto.txt \
  -out mensaje_cifrado.bin \
  -pass file:RedKey.txt
\end{lstlisting}

\textbf{Explicación de parámetros:}
\begin{itemize}
    \item \texttt{-aes-256-cbc}: Algoritmo de cifrado robusto.
    \item \texttt{-salt -pbkdf2}: Añade seguridad contra ataques de diccionario.
    \item \texttt{-pass file:RedKey.txt}: Le dice a openssl que la contraseña está dentro del archivo \texttt{RedKey.txt}.
\end{itemize}

Ahora \texttt{mensaje\_cifrado.bin} es un archivo ilegible sin la llave.

\subsection{Paso 3 --- Ocultar la llave en RedTeam.jpg}

\begin{lstlisting}[language=bash, caption={Ocultar llave con Steghide}]
steghide embed -cf RedTeam.jpg -ef RedKey.txt -sf RedTeamKey.jpg
\end{lstlisting}

\begin{notebox}[Nota importante]
Te pedirá una contraseña de steghide. Anótala, la necesitará el receptor para abrir esta imagen. Ejemplo: \texttt{pass\_red}.
\end{notebox}

\subsection{Paso 4 --- Ocultar el archivo cifrado en BlueTeam.jpg}

\begin{lstlisting}[language=bash, caption={Ocultar mensaje cifrado con Steghide}]
steghide embed -cf BlueTeam.jpg -ef mensaje_cifrado.bin -sf BlueTeamFinal.jpg
\end{lstlisting}

\begin{notebox}[Nota importante]
Te pedirá otra contraseña de steghide. Ejemplo: \texttt{pass\_blue}.
\end{notebox}

\subsection{Paso 5 --- Limpiar archivos temporales}

\begin{lstlisting}[language=bash, caption={Eliminar archivos sensibles}]
shred -u RedKey.txt mensaje_secreto.txt mensaje_cifrado.bin
\end{lstlisting}

\textbf{Resultado final:} Solo te quedan \texttt{RedTeamKey.jpg} y \texttt{BlueTeamFinal.jpg} para enviar al receptor.

\section{Fase 2: Extracción (Receptor)}

\subsection{Paso 1 --- Extraer la llave de RedTeam}

\begin{lstlisting}[language=bash, caption={Extraer llave de RedTeam}]
steghide extract -sf RedTeamKey.jpg -xf llave_recuperada.txt
\end{lstlisting}

Introduce la contraseña de steghide de RedTeam (\texttt{pass\_red}).

\subsection{Paso 2 --- Extraer el archivo cifrado de BlueTeam}

\begin{lstlisting}[language=bash, caption={Extraer mensaje cifrado de BlueTeam}]
steghide extract -sf BlueTeamFinal.jpg -xf mensaje_cifrado_recuperado.bin
\end{lstlisting}

Introduce la contraseña de steghide de BlueTeam (\texttt{pass\_blue}).

\subsection{Paso 3 --- Descifrar el mensaje usando la llave recuperada}

\begin{lstlisting}[language=bash, caption={Descifrar mensaje con OpenSSL}]
openssl enc -aes-256-cbc -d -pbkdf2 \
  -in mensaje_cifrado_recuperado.bin \
  -out mensaje_final.txt \
  -pass file:llave_recuperada.txt
\end{lstlisting}

\textbf{Explicación:}
\begin{itemize}
    \item \texttt{-d}: Modo \textbf{descifrado} (decrypt).
    \item \texttt{-pass file:llave\_recuperada.txt}: Usa la llave que extrajo de RedTeam.
\end{itemize}

\subsection{Paso 4 --- Leer el mensaje}

\begin{lstlisting}[language=bash, caption={Visualizar mensaje final}]
cat mensaje_final.txt
\end{lstlisting}

\textbf{Salida esperada:}
\begin{verbatim}
El punto de encuentro es el muelle 7 a las 03:00 AM. Destruye este mensaje.
\end{verbatim}

\section{Resumen Visual del Flujo}

\begin{figure}[h]
\centering
\begin{tikzpicture}[
    node distance=1.5cm,
    box/.style={rectangle, draw=blue!75!black, fill=blue!5, thick, minimum width=3cm, minimum height=0.8cm, align=center},
    arrow/.style={->, >=stealth, thick, color=blue!75!black}
]

% Emisor
\node[box] (llave) {RedKey.txt};
\node[box, right=of llave] (redteam) {RedTeam.jpg};
\node[box, right=of redteam] (redfinal) {RedTeamKey.jpg};

\node[box, below=1.5cm of llave] (mensaje) {mensaje\_secreto.txt};
\node[box, right=of mensaje] (cifrado) {mensaje\_cifrado.bin};
\node[box, below=of redfinal] (blueteam) {BlueTeam.jpg};
\node[box, right=of blueteam] (bluefinal) {BlueTeamFinal.jpg};

% Flechas emisor
\draw[arrow] (llave) -- node[above] {steghide} (redteam);
\draw[arrow] (redteam) -- node[above] {embed} (redfinal);

\draw[arrow] (llave) |- node[pos=0.25, left] {openssl cifrar} (mensaje);
\draw[arrow] (mensaje) -- node[above] {genera} (cifrado);
\draw[arrow] (cifrado) |- node[pos=0.75, right] {steghide} (blueteam);
\draw[arrow] (blueteam) -- node[above] {embed} (bluefinal);

% Receptor
\node[box, below=2cm of redfinal] (redfinal2) {RedTeamKey.jpg};
\node[box, right=of redfinal2] (llave2) {llave\_recuperada.txt};

\node[box, below=of bluefinal] (bluefinal2) {BlueTeamFinal.jpg};
\node[box, right=of bluefinal2] (cifrado2) {mensaje\_cifrado.bin};
\node[box, below=of llave2] (final) {mensaje\_final.txt};

% Flechas receptor
\draw[arrow] (redfinal2) -- node[above] {steghide extract} (llave2);
\draw[arrow] (bluefinal2) -- node[above] {steghide extract} (cifrado2);
\draw[arrow] (llave2) |- node[pos=0.25, left] {openssl descifrar} (cifrado2);
\draw[arrow] (cifrado2) -- node[above] {genera} (final);

% Etiquetas
\node[above=0.3cm of redfinal, font=\bfseries, color=red!75!black] {EMISOR};
\node[below=0.3cm of final, font=\bfseries, color=green!50!black] {RECEPTOR};

\end{tikzpicture}
\caption{Flujo completo de cifrado y descifrado}
\end{figure}

\section{Notas Importantes}

\begin{enumerate}
    \item \textbf{Las contraseñas de steghide} (\texttt{pass\_red} y \texttt{pass\_blue}) deben comunicarse al receptor por un \textbf{canal separado y seguro} (por ejemplo, de voz, Signal, o en persona). Si las mandas por el mismo chat que las imágenes, pierdes toda la seguridad.

    \item \textbf{Si la llave se corrompe} (aunque sea un solo carácter), \texttt{openssl} no podrá descifrar y dará un error de \textit{"bad decrypt"}. Por eso es buena idea que la llave sea algo fácil de verificar visualmente.

    \item \textbf{Las imágenes originales} (\texttt{RedTeam.jpg} y \texttt{BlueTeam.jpg}) no se modifican. Steghide crea copias nuevas (\texttt{RedTeamKey.jpg} y \texttt{BlueTeamFinal.jpg}). No envíes las originales.

    \item \textbf{El orden importa}: El receptor \textit{debe} extraer primero RedTeam para obtener la llave antes de poder abrir BlueTeam. Es buena idea indicarle el orden de alguna forma (el nombre ya lo sugiere: rojo primero, azul después).

    \item \textbf{Compatibilidad de formatos}: Steghide soporta formatos de imagen como \texttt{.jpg}, \texttt{.bmp}, \texttt{.wav} y \texttt{.au}. Asegúrate de usar uno compatible.

    \item \textbf{Seguridad adicional}: Considera usar \texttt{shred} o herramientas similares para eliminar de forma segura los archivos temporales sensibles.
\end{enumerate}

\section{Comandos Útiles Adicionales}

\subsection{Ver información de una imagen con datos ocultos}
\begin{lstlisting}[language=bash]
steghide info RedTeamKey.jpg
\end{lstlisting}

Esto te dirá si la imagen tiene datos ocultos, qué formato tiene el archivo oculto y qué algoritmo de cifrado se usó, pero \textbf{no} te mostrará el contenido ni te dejará extraerlo sin la contraseña correcta.

\subsection{Parámetros de Steghide}

\begin{table}[h]
\centering
\begin{tabular}{|l|l|l|}
\hline
\textbf{Parámetro} & \textbf{Significado} & \textbf{Uso} \\
\hline
\texttt{-cf} & Cover File & Archivo de portada (imagen original) \\
\hline
\texttt{-ef} & Embed File & Archivo que quieres ocultar \\
\hline
\texttt{-sf} & Stego File & Archivo de salida (imagen con datos ocultos) \\
\hline
\texttt{-xf} & eXtract File & Archivo de salida al extraer (opcional) \\
\hline
\texttt{-p} & Passphrase & Contraseña (opcional) \\
\hline
\end{tabular}
\end{table}

\end{document}